% framedschemapage_example.tex
% Demonstrates the \framedschemapage command for full-page rotated diagrams
% Version: 7.2.0

\documentclass[book]{hebrew-academic-template}

% Bibliography
\addbibresource{example_references.bib}

% Document metadata
\hebrewtitle{דוגמה לפקודת framedschemapage}
\englishtitle{Framed Schema Page Example}
\hebrewauthor{ד"ר סגל יורם}

\begin{document}

\maketitle
\tableofcontents

% =============================================================================
\hebrewchapter{מבוא}
% =============================================================================

פרק זה מדגים את השימוש בפקודה \en{\texttt{\textbackslash framedschemapage}} שנוספה בגרסה \en{7.2.0} של התבנית.

\hebrewsection{מטרת הפקודה}

הפקודה \en{\texttt{\textbackslash framedschemapage}} מיועדת להצגת דיאגרמות גדולות בעמוד מלא:
\begin{itemize}
    \item דיאגרמות סכמה של בסיס נתונים
    \item דיאגרמות ארכיטקטורה
    \item תרשימי זרימה
    \item כל תמונה רחבה שיש להציג בכיוון לאורך
\end{itemize}

\hebrewsection{תחביר הפקודה}

\begin{english}
\begin{verbatim}
% Basic usage (default size 0.75\textwidth):
\framedschemapage{path/to/image.png}

% Custom size:
\framedschemapage[0.8\textwidth]{path/to/image.png}

% Smaller size:
\framedschemapage[0.6\textwidth]{path/to/image.png}
\end{verbatim}
\end{english}

\hebrewsection{תכונות הפקודה}

\begin{enumerate}
    \item \textbf{סיבוב אוטומטי} --- הפקודה מסובבת את התמונה ב-\en{90°} נגד כיוון השעון
    \item \textbf{מיקום מרכזי} --- התמונה ממוקמת במרכז העמוד
    \item \textbf{מסגרת דקורטיבית} --- מסגרת פשוטה ואלגנטית סביב העמוד
    \item \textbf{תמיכה ב-\en{RTL}} --- עובדת נכון במסמכים בעברית
    \item \textbf{עמוד ריק} --- ללא כותרות עליונות ותחתונות
\end{enumerate}

% =============================================================================
\hebrewchapter{דוגמה לשימוש}
% =============================================================================

בעמוד הבא מוצגת דיאגרמת ארכיטקטורה באמצעות הפקודה \en{\texttt{\textbackslash framedschemapage}}.

% Insert the framed schema page
\framedschemapage{images/fig_architecture.png}

% =============================================================================
\hebrewchapter{פרטים טכניים}
% =============================================================================

\hebrewsection{איך זה עובד}

הפקודה משתמשת בטכניקות הבאות:
\begin{itemize}
    \item \en{\texttt{\textbackslash clearpage}} --- להבטחת עמוד חדש
    \item \en{\texttt{\textbackslash thispagestyle\{empty\}}} --- להסרת כותרות
    \item \en{\texttt{english} environment} --- לטיפול נכון בכיווניות
    \item \en{\texttt{\textbackslash rotatebox\{-90\}}} --- לסיבוב התמונה
    \item \en{TikZ overlay} --- לציור המסגרת
\end{itemize}

\hebrewsection{דוגמה עם גודל מותאם}

ניתן לשנות את גודל התמונה באמצעות הפרמטר האופציונלי:

\begin{english}
\begin{verbatim}
% Larger image (80% of text width):
\framedschemapage[0.8\textwidth]{images/schema.png}

% Smaller image (60% of text width):
\framedschemapage[0.6\textwidth]{images/schema.png}
\end{verbatim}
\end{english}

% Another example with different size
\framedschemapage[0.65\textwidth]{images/fig_data.png}

% =============================================================================
\hebrewchapter{סיכום}
% =============================================================================

הפקודה \en{\texttt{\textbackslash framedschemapage}} מספקת פתרון פשוט ואלגנטי להצגת דיאגרמות גדולות בעמוד מלא, עם תמיכה מלאה במסמכים בעברית.

\end{document}
